\chapter{واجهات نظام التشغيل}
\label{CH:UNIX}

تتمثل مهمة نظام التشغيل (Operating System، OS) في مشاركة حاسوب واحد بين عدة برامج، وتقديم مجموعة من الخدمات المحسّنة والمفيدة التي تتجاوز ما يوفره العتاد الصلب بمفرده. يحيط نظام التشغيل المعالج والذاكرة وأجهزة الإدخال والإخراج المادية بتجريب برمجية رفيعة المستوى، مما يضمن التوزيع العادل للموارد وتوفير بيئة تنفيذ معزولة وآمنة لكل برنامج.

يوفر نظام التشغيل خدماته للبرامج التطبيقية الخاصة بالمستخدم عبر واجهة برمجة محددة (Interface). وقد أثبتت التجارب أن تصميم واجهة نظام تشغيل جيدة ومتوازنة يُعد تحديًا هندسيًا كبيرًا؛ إذ يجب أن تكون الواجهة بسيطة وأنيقة ومحدودة العدد من جهة، ومكفولة القدرة والتعبير عن العمليات المعقدة من جهة أخرى. تميل أنظمة التشغيل إلى تجميع هذه الخدمات داخل نواة مركزية تحمي نفسها والعتاد من البرامج الخبيثة أو المعطوبة.

يستخدم هذا الكتاب نظام تشغيل حقيقيًا ومكتمل الأركان كمثال عملي لتوضيح مفاهيم تصميم أنظمة التشغيل، وهو نظام \texttt{xv6}. يوفر \texttt{xv6} الواجهات والخدمات الأساسية التي قدمها نظام Unix التاريخي الذي صممه كين طومسون ودنيس ريتشي، ولكنه أُعيدت كودته بالكامل بلغة C الحديثة ليعمل على المعالجات المعتمدة على معمارية \textsf{RISC-V} 64-bit متعددة الأنوية (Multicore CPUs).

كما يوضح الشكل~\ref{fig:os}، يأخذ \texttt{xv6} الشكل التقليدي المتمثل في وجود \textbf{النواة} (Kernel)، وهي برنامج ممتاز يشتغل في وضع صلاحيات ممتاز ويقدم الخدمات للبرامج الجارية الأخرى. يُطلق على كل برنامج جاري ينفذ التعليمات اسم \textbf{العملية} (Process)، وتملك كل عملية فضاء عناوين افتراضيًا (Virtual Address Space) يحتوي على الذاكرة الخاصة بها (التعليمات، البيانات، والمكدس)، بالإضافة إلى حالة خاصة بالنواة لتعقب سير تنفيذها.

عندما تحتاج العملية إلى استدعاء خدمة من خدمات النواة، فإنها تنفذ \textbf{استدعاء نظام} (System Call)، وهو أحد الاستدعاءات المعرفة في واجهة نظام التشغيل. تنتقل السيطرة فور تنفيذ استدعاء النظام من وضع المستخدم (User Mode) غير الممتاز إلى وضع المشرف (Supervisor Mode) الممتاز داخل النواة، فتتفحص النواة وسائط الطلب، وتنفذ الخدمة، ثم تعيد النتيجة والسيطرة إلى برنامج المستخدم.

وكما سنشرح بالتفصيل في الفصول القادمة، تستخدم النواة آليات الحماية العتادية التي يوفرها المعالج لحظر برامج المستخدم من الوصول المباشر للسجلات العتادية الحساسة أو ذاكرة العمليات الأخرى. تعتمد النواة على هذه الحماية لتحقيق *العزل القوي* (Strong Isolation) بين العمليات ومنع عملية واحدة معطوبة من شل حركة النظام بأكمله.

تُشكل مجموعة استدعاءات النظام التي توفرها النواة الواجهة الفعلية التي تراها برامج المستخدم. يوفر نظام \texttt{xv6} مجموعة فرعية مصغرة ولكنها مكتملة من خدمات واستدعاءات نظام Unix القياسية. يوضح الجدول~\ref{tab:syscalls} كافة استدعاءات النظام المعتمدة في \texttt{xv6}.

استعرض بقية هذا الفصل الخدمات الأساسية التي يقدمها \texttt{xv6}---العمليات، الذاكرة، واصفات الملفات، الأنابيب، ونظام الملفات---ويوضحها بالعديد من الأمثلة والأكواد البرمجية المقتطعة من مغلّف الأوامر (Shell).

تُعد مغلّف الأوامر (\texttt{sh}) في نظام Unix برنامجًا عاديًا بلغة المستخدم يقرأ الأوامر من المستخدم وينفذها، وليست جزءًا ثابتا من كود النواة. توضح هذه الحقيقة مرونة واجهة Unix؛ إذ يمكن للمستخدمين استبدال مغلّف الأوامر الافتراضي بأي مغلّف آخر. يمكن العثور على تنفيذ مغلّف الأوامر في \texttt{xv6} داخل الملف \texttt{user/sh.c}.

\section{العمليات والذاكرة}
\label{sec:process_memory}

تتكون عملية \texttt{xv6} من ذاكرة فضاء المستخدم (المشتملة على التعليمات التنفيذية والبيانات والمكدس) وحالة خاصة بالعملية خافية ومحفوظة داخل النواة. يطبق \texttt{xv6} تقنية مشاركة الوقت (Time-sharing) بين العمليات؛ إذ يوزع أنوية المعالج المتاحة بين العمليات الجاهزة للتنفيذ، وعندما تنتظر عملية حدوث خطوة إدخال أو إخراج، يوقفها النواة ويعين المعالج لعملية أخرى.

عندما لا تكون العملية قيد التنفيذ، تحفظ النواة سجلات المعالج الخاصة بها وتسترجعها عند إعادة إشغال المعالج بها. تُحدد كل عملية بواسطة عدد صحيح موجب يُدعى \textbf{معرّف العملية} (\texttt{PID} - Process Identifier).

يمكن للعملية إنشاء عملية جديدة باستخدام استدعاء النظام \texttt{fork()}. ينشئ الاستدعاء \texttt{fork()} عملية جديدة تُدعى \textbf{العملية الابن} (Child Process)، وتكون نسخة مطابقة تمامًا للعملية المداعية التي تُدعى \textbf{العملية الأب} (Parent Process)، حيث تحصل العملية الابن على نفس محتويات ذاكرة الأب ونفس السجلات ونفس واصفات الملفات المفتوحة.

يرجع الاستدعاء \texttt{fork()} بقيمة عدادية في كلتا العمليتين:
\begin{itemize}
  \item في العملية الأب: يرجع \texttt{fork()} معرّف العملية الابن (\texttt{PID}).
  \item في العملية الابن: يرجع \texttt{fork()} القيمة \texttt{0}.
\end{itemize}
تسمح هذه الآلية الأنيقة لكلا العمليتين بالتفرغ والتمييز بين كود الأب وكود الابن باستخدام جملة شرطية بسيطة مثل \texttt{if (fork() == 0)}.

ينفذ استدعاء النظام \texttt{exit(status)} إيقاف تنفيذ العملية الحالية وتدميرها، مع التمرير برقم حالة الخروج (\texttt{status}) إلى النواة. وتستخدم العملية الأب استدعاء النظام \texttt{wait(\&status)} للانتظار حتى تنهي إحدى عملياتها الابن تنفيذها؛ حيث يرجع \texttt{wait()} معرّف الابن المنتهي ويخزن رقم حالة الخروج في المتغير الممرر.

تحمل دالة استدعاء النظام \texttt{exec(filename, argv)} ملفًا تنفيذيًا من نظام الملفات بتنسيق ELF، وتحمله داخل فضاء عناوين العملية الحالية بدلاً من الذاكرة القديمة، ثم تبدأ بتنفيذ نقطة الدخول المحددة في الملف الممرر مع تمرير الوسائط \texttt{argv}.

توفر واجهة \texttt{xv6} استدعاء النظام \texttt{sbrk(n)} لتوسيع أو تقليص الذاكرة المخصصة للعملية بمقدار \texttt{n} بايت، حيث تعيد النواة عنوان بداية الذاكرة المخصصة حديثًا.

\section{الإدخال والإخراج وواصفات الملفات}
\label{sec:io_fds}

\textbf{واصف الملف} (File Descriptor) هو عدد صحيح صغير يشير إلى مورد مفتوح يُدار بواسطة النواة ويمكن للعملية القراءة منه أو الكتابة فيه. تشمل الموارد المفتوحة الملفات العادية، الدلائل، الأنابيب، وأجهزة الإدخال والإخراج المباشرة مثل شاشة المنصة (Console).

تُعين اتفاقية Unix القياسية أول ثلاثة واصفات ملفات لكل عملية فور إنشائها:
\begin{itemize}
  \item الواصف \texttt{0}: \textbf{الإدخال القياسي} (Standard Input - \texttt{stdin}).
  \item الواصف \texttt{1}: \textbf{الخرج القياسي} (Standard Output - \texttt{stdout}).
  \item الواصف \texttt{2}: \textbf{الخطأ القياسي} (Standard Error - \texttt{stderr}).
\end{itemize}

يُنفذ استدعاء النظام \texttt{read(fd, buf, n)} قراءة ما يصل إلى \texttt{n} بايت من واصف الملف \texttt{fd} ونسخها في ذاكرة المؤشر \texttt{buf}، ويرجع عدد البايتات المأخوذة فعليًا، أو \texttt{0} عند الوصول لنهاية الملف (EOF). بينما ينفذ \texttt{write(fd, buf, n)} كتابة \texttt{n} بايت من \texttt{buf} إلى \texttt{fd}.

يستدعي فتح الملفات الاستدعاء \texttt{open(filename, flags)}، حيث يحدد \texttt{flags} نمط الفتح (مثل \texttt{O\_RDONLY} للقراءة فقط، \texttt{O\_WRONLY} للكتابة، \texttt{O\_CREATE} لإنشاء الملف إذا لم يكن موجودًا). تعيد النواة أصغر واصف ملف غير مستخدم في العملية الحالية.

يوفر الاستدعاء \texttt{close(fd)} إغلاق واصف الملف وتحريره ليكون متاحًا للاستخدام مستقبلاً. بينما يقوم الاستدعاء \texttt{dup(fd)} بإنشاء واصف ملف جديد يشير إلى نفس المورد المفتوح ويشارك نفس مؤشر المواضع المحدث.

\section{الأنابيب}
\label{sec:pipes}

\textbf{الأنبوب} (Pipe) هو قناة اتصل أحادية الاتجاه تُدار داخل النواة لربط عملية كاتب بعملية قارئ. يوفر الأنبوب مفرغًا مؤقتًا في الذاكرة النواتية يُمكن عملية من كتابة تدفق بيانات وتقوم العملية الأخرى بقراءتها بنفس الترتيب.

ينشئ استدعاء النظام \texttt{pipe(p)} أنبوبًا جديدًا ويخزن واصفين اثنين في المصفوفة \texttt{p}:
\begin{itemize}
  \item \texttt{p[0]}: واصف الملف المخصص للقراءة من الأنبوب.
  \item \texttt{p[1]}: واصف الملف المخصص للكتابة في الأنبوب.
\end{itemize}

إذا حاولت عملية القراءة من أنبوب فارغ لا تحتوي ذاكرته المؤقتة على بايتات، يحجب (\textbf{blocking}) النظام العملية القارئة وتدخل في حالة تعليق حتى تقوم العملية الكاتبة بضخ بيانات جديدة أو إغلاق واصف الكتابة.

تستغل مغلّفات الأوامر الأنابيب لإنشاء سلاسل الأوامر المترابطة (Pipelines) مثل \texttt{grep fork sh.c | wc -l}، حيث يُحول الخرج القياسي للأمر الأول مباشرةً ليكون هو الإدخال القياسي للأمر الثاني دون الحاجة إلى تخزين البيانات في ملفات وسيطة على القرص.

\section{نظام الملفات}
\label{sec:filesystem}

يوفر نظام الملفات في \texttt{xv6} التسمية والتخزين الدائم للملفات والدلائل على القرص الصلب. يتكون نظام الملفات من شجرة ممتدة من \textbf{الدلائل} (Directories) والملفات العادية المحددة بواسطة مسارات نصية المظهر (Pathnames).

يبدأ الهيكل الشجري من الدليل الجذر المعبر عنه بالرمز \texttt{/}، ويتكون كل دليل من قائمة مدخلات تربط الأسماء النصية بـ \textbf{أرقام العقد الفهرسية} (Inode Numbers).

يوفر \texttt{xv6} مجموعة من استدعاءات النظام لإدارة نظام الملفات:
\begin{itemize}
  \item \texttt{mkdir(dirname)}: إنشاء دليل جديد.
  \item \texttt{mknod(path, major, minor)}: إنشاء ملف جهاز عتادي خاص.
  \item \texttt{fstat(fd, \&st)}: جلب البيانات الوصفية (Metadata) للملف المفتوح وتعبئتها في الهيكل \texttt{struct stat}.
  \item \texttt{link(oldpath, newpath)}: إنشاء وصلة صلبة (Hard Link) تشير إلى نفس العقدة الفهرسية.
  \item \texttt{unlink(path)}: حذف الاسم من الدليل وإلغاء ربطه بالعقدة الفهرسية، وعندما يصل عدد الوصلات إلى \texttt{0} وتُغلق جميع الواصفات، تقوم النواة بتحرير كتل الملف ماديًا من القرص.
  \item \texttt{chdir(path)}: تغيير الدليل الحالي للعملية.
\end{itemize}

\section{العالم الحقيقي}
\label{sec:real_world}

تتبع أنظمة التشغيل الحديثة المعتمدة على معيار POSIX (مثل Linux و FreeBSD و macOS) نفس المفاهيم الأساسية لواجهات Unix التي يقدمها \texttt{xv6}. ومع ذلك، توفر الأنظمة التجارية العديد من استدعاءات النظام الإضافية لتحسين الأداء والأمان والتزامن المتقدم (مثل \texttt{select} و \texttt{poll} و \texttt{epoll} للتنقل غير الحاجب - non-blocking I/O).

كما تدعم هذه الأنظمة العديد من أنواع أنظمة الملفات المتقدمة (مثل ext4، ZFS، و Btrfs) التي توفر التخزين المؤقت الذكي، والمساحات الكبيرة، والصلاحيات التفصيلية عبر قوائم التحكم بالوصول (ACLs).

\section{تمارين}
\label{sec:exercises}

\begin{enumerate}
  \item اكتب برنامجًا بلغة C لـ \texttt{xv6} يستخدم الأنابيب لنقل مصفوفة بايتات بين عمليتين أصل وابن مع قياس الأداء.
  \item قم بتعديل مغلّف الأوامر \texttt{user/sh.c} لدعم إعادة توجيه الإدخال والخرج معًا في نفس السطر.
  \item أضف خيارًا لإنشاء وصلات رمزية (Symbolic Links) واشرح كيف تختلف عن الوصلات الصلبة (Hard Links).
\end{enumerate}
